\documentclass[conference]{IEEEtran}
\IEEEoverridecommandlockouts
% The preceding line is only needed to identify funding in the first footnote. If that is unneeded, please comment it out.
\usepackage{cite}
\usepackage{amsmath,amssymb,amsfonts}
\usepackage{algorithmic}
\usepackage{graphicx}
\usepackage{textcomp}
\usepackage{xcolor}
\usepackage{url}
\usepackage[hidelinks]{hyperref}
\def\BibTeX{{\rm B\kern-.05em{\sc i\kern-.025em b}\kern-.08em
    T\kern-.1667em\lower.7ex\hbox{E}\kern-.125emX}}
\begin{document}

\title{Data Science and Machine Learning Project Article\\
{\footnotesize \textsuperscript{*}Note: Sub-titles are not captured in Xplore and
should not be used}
\thanks{Identify applicable funding agency here. If none, delete this.}
}

\author{\IEEEauthorblockN{1\textsuperscript{st} Md Rajaul Karim}
\IEEEauthorblockA{\textit{Dept. of Information Technology} \\
\textit{Georgia Southern University}\\
Statesboro, Georgia, USA \\
ORCID: 0000-0003-4152-7290}}

% \and
% \IEEEauthorblockN{2\textsuperscript{nd} Vijayalakshmi Ramasamy}
% \IEEEauthorblockA{\textit{Dept. of Computer Science} \\
% \textit{Georgia Southern University}\\
% Statesboro, Georgia, USA \\
% email address or ORCID} }


% \and
% \IEEEauthorblockN{3\textsuperscript{rd} Given Name Surname}
% \IEEEauthorblockA{\textit{dept. name of organization (of Aff.)} \\
% \textit{name of organization (of Aff.)}\\
% City, Country \\
% email address or ORCID}
% \and
% \IEEEauthorblockN{4\textsuperscript{th} Given Name Surname}
% \IEEEauthorblockA{\textit{dept. name of organization (of Aff.)} \\
% \textit{name of organization (of Aff.)}\\
% City, Country \\
% email address or ORCID}
% \and
% \IEEEauthorblockN{5\textsuperscript{th} Given Name Surname}
% \IEEEauthorblockA{\textit{dept. name of organization (of Aff.)} \\
% \textit{name of organization (of Aff.)}\\
% City, Country \\
% email address or ORCID}
% \and
% \IEEEauthorblockN{6\textsuperscript{th} Given Name Surname}
% \IEEEauthorblockA{\textit{dept. name of organization (of Aff.)} \\
% \textit{name of organization (of Aff.)}\\
% City, Country \\
% email address or ORCID}
% }

\maketitle

\begin{abstract}
This paper presents an overview of a data science and machine learning project. The objective is to demonstrate the creation of a research article using LaTeX and the IEEE conference template. This work also emphasizes structuring sections and integrating external resources such as a GitHub repository. This abstract summarizes the purpose of the paper in a few sentences.

\end{abstract} 


\begin{IEEEkeywords}
LaTeX, IEEE conference format, data science, machine learning, academic writing
\end{IEEEkeywords}

\noindent\textbf{GitHub Repository:} 
\color{blue}\url{https://github.com/mk19409-prog/DS_and_ML_Project_Article_LaTex}\color{black}





\section{Introduction}
This section introduces the problem domain and the motivation behind the project. The section introduces the topic and explains why it is important.

\section{Literature Review}
This section summarizes existing research or tools related to the project topic including referencing article \cite{karim2025comprehensive}.

\section{Main Methodology}
This section describes the dataset, tools, and methods used in the project. Describes the methods or steps taken to complete the project.

For this assignment, the project methodology includes using LaTeX for document preparation, following the IEEE template, and organizing content into structured sections.
All materials are maintained in a \color{blue}\href{https://github.com/mk19409-prog/DS_and_ML_Project_Article_LaTex}{GitHub repository: DS and ML Project} \color{black}
.


\section{Result Analysis}
This section highlights preliminary results and discusses observations. This basically discusses what we achieved or produced through the research.

\section{Discussion and Future Work}
This section discusses the findings of the project and highlights key observations from the analysis. Reflects on what was learned and what could be done next. It demonstrates critical thinking, noting the strengths of the work and suggesting improvements for future assignments or projects.

\section{Conclusion}
This section summarizes the main points of the paper and what was accomplished. In this assignment, it demonstrates that we successfully used LaTeX, followed the IEEE template, and integrated a GitHub repository, rather than reporting scientific results. It highlights the completion of the document with proper formatting.


\section*{Acknowledgment}
The author thanks the instructor and department for guidance, funding, and support.


\bibliographystyle{IEEEtran}

\bibliography{reference.bib}


\end{document}
